\chapter{Conclusion}
\label{chapter:conclusion}

In this thesis, we studied finite-state devices whose graphs contain a separator, and
showed how a separator can be used to obtain deterministic representations that are
smaller than those produced by the classical constructions.

\section{Summary}

\paragraph{Quasi-determinization.}
A $k$-separated quasi-deterministic finite-state automaton keeps one deterministic
quasi-state per component and passes information between the components only through
the separator. Every $k$-separated finite-state automaton has an equivalent
quasi-deterministic finite-state automaton whose $i$-th component has at most $2^{|Q_i|
+ |X|}$ quasi-states, whereas a deterministic finite-state automaton may need up to
$2^{|Q|}$ states. Its correctness rests on \Cref{lem:quasidet_invariant}: the
quasi-states of a reachable state are exactly the projections $Reach(w) \cap
\mathcal{C}(Q_i)$ of the corresponding state of the subset construction. In particular,
every reachable separator state belongs to all quasi-states, and the separator transition
function is what keeps the components synchronized. Words are processed in
$\mathcal{O}(k)$ time per symbol (\Cref{alg:qdfsa_traversal}).

\paragraph{Construction time.}
The small size of a quasi-deterministic finite-state automaton does not imply that it can
be constructed quickly. \Cref{alg:quasi_det_sfsa} explores the reachable states of
$\mathcal{A}'$, and by \Cref{prop:quasidet_iterations} there are exactly as many of them
as states of the deterministic finite-state automaton constructed by
\Cref{thm:rabin_scott}. The finite-state automata $\mathcal{A}_k$ show that this number
can be exponentially larger than the sizes of the components. Since the tables of a
quasi-deterministic finite-state automaton are indexed by component quasi-states and
separator states only, they can instead be computed for each component separately, by
passing all subsets of the separator to the extension functions.
\Cref{alg:quasi_det_sfsa_components} does this in $\mathcal{O}((|\Sigma| + 2^{|X|}) |Q|
2^{|X|} \sum_{i=1}^k 2^{|Q_i|})$ time (\Cref{thm:quasidet_components}), independently of
the size of the deterministic finite-state automaton. The price is that the constructed
device may contain quasi-states that are not components of reachable states. All
quasi-states are $\epsilon$-closed and determined by their essential states and their
separator states, so $|Q_i|$ can be replaced by the number $m_i$ of essential states of
$Q_i$ in both bounds (\Cref{lem:quasidet_essential,thm:quasidet_components_essential}).

\paragraph{Regular expressions.}
Finite-state automata constructed from regular expressions by the algorithm described in
\Cref{thm:thompson} have separators of size $2$ whose components have at most
$\frac{2}{3}|Q|$ states, and \Cref{alg:thompson_separator} finds them in linear time
(\Cref{thm:thompson_separator}). For quasi-determinization, the relevant measure is the
number $m$ of symbol occurrences: the subset construction produces at most $2^{m+1}$
states, while $|Q| \geq 2m$. The symbol-balanced variant of
\Cref{alg:thompson_separator} finds separators of size $2$ whose components have at most
$\frac{2}{3}m + 1$ essential states (\Cref{thm:thompson_separator_symbols}). Together
with \Cref{alg:quasi_det_sfsa_components}, every regular expression can be
quasi-determinized in $\mathcal{O}(|\Sigma| |Q| 2^{\frac{2}{3}m})$ time
(\Cref{cor:quasidet_regex}). In the experiments of
\Cref{tbl:state_comparison,tbl:time_comparison}, for the expression
$(0|1)*.0.(0|1)^{20}$ the deterministic finite-state automaton has $2097153$ states.
\Cref{alg:quasi_det_sfsa_components} constructs an equivalent $3$-separated
quasi-deterministic finite-state automaton with $7282$ quasi-states in $7281$ iterations,
where \Cref{alg:quasi_det_sfsa} needs $2097153$ iterations.

\paragraph{Boundaries.}
A component shares transitions only with some of the separator states, its boundary $X_i$.
Recording only these in its quasi-states, while the separator state is still passed as a
subset of $X$, gives an equivalent quasi-deterministic finite-state automaton with at most
$2^{m_i + |X_i|}$ quasi-states in the $i$-th component, constructed in
$\mathcal{O}(\sum_{i=1}^k (|\Sigma| + 2^{|X_i|}) |Q| 2^{m_i + |X_i|})$ time
(\Cref{thm:quasidet_boundaries}). For regular expressions, the boundaries stay small when
the number of components grows. For $(0|1)*.0.(0|1)^{20}$ and $16$ components, the
construction produces $185$ quasi-states (\Cref{tbl:boundary_comparison}).

\paragraph{Transducers.}
The same approach applies to trimmed functional $k$-separated finite-state transducers
whose function has bounded variation. Here, quasi-states carry pending outputs, and the
outputs of the components have to be normalized by their common prefix. This prefix
depends on all components, so it cannot be computed by the quasi-state transition
functions $\tau_i$ alone. The construction of
\Cref{section:quasi_subsequentialization} therefore passes it through the separator
states, and the extension functions $\eta_i$ perform the normalization.
\Cref{lem:quasisubseq_invariant} characterizes the reachable states, and bounded
variation guarantees that only finitely many of them exist
(\Cref{lem:quasisubseq_completion}).

\section{Formal verification}

The main results of \Cref{chapter:separators_finite_state_automata,chapter:separators_finite_state_transducers}
are formalized in the Lean theorem prover, using the Mathlib library, in the directory
\texttt{lean} of the accompanying repository. The formalization covers the
quasi-determinization and quasi-subsequentialization constructions and their
correctness, the completion lemmas, the equivalence of quasi-deterministic and
deterministic finite-state automata, the complexity bounds of the traversal algorithms
and of \Cref{alg:quasi_det_sfsa_components}, including the refined bounds in terms of
essential states (\Cref{lem:quasidet_essential,thm:quasidet_components_essential}) and
the construction with boundaries (\Cref{lem:boundary_invariant,lem:boundary_essential,thm:quasidet_boundaries}), the
iteration count of
\Cref{alg:quasi_det_sfsa} together with the finite-state automata $\mathcal{A}_k$, and the
separators of \Cref{alg:thompson_separator} and of its symbol-balanced variant together with
\Cref{cor:quasidet_regex}. It does not cover \Cref{prop:essential_dfa}, the planarity of the
graphs of finite-state automata constructed from regular expressions, the results cited
from the literature (\Cref{thm:rabin_scott,thm:thompson,thm:choffrut_mohri}), and the
correctness of the pseudocode of the construction algorithms line by line. The
algorithms are modelled by the sequences of states they process.

\section{Future work}

\paragraph{Per-component quasi-subsequentialization.}
\Cref{alg:quasi_det_sfsa_components} relies on the finiteness of the set of separator
states $2^X$. For transducers, separator states also carry pending outputs and a common
prefix, so there are infinitely many candidates. A per-component construction would
need a bound on the pending outputs, such as the one derived from bounded variation in
the proof of \Cref{lem:quasisubseq_completion}, and a way to enumerate only the relevant
separator states.

\paragraph{Experiments with transducers.}
The implementation covers finite-state automata only. Implementing the
quasi-subsequentialization of \Cref{section:quasi_subsequentialization} would allow the
same comparison for transducers, for example for transducers constructed from rewrite
rules.

\paragraph{Size of the constructed devices.}
\Cref{alg:quasi_det_sfsa_components} trades construction time for size: it may tabulate
quasi-states that no reachable state uses. These could be removed afterwards by a
traversal of the reachable states, but such a traversal again processes as many states as
the deterministic finite-state automaton has. Bounding the number of unused quasi-states,
or avoiding them without enumerating global states, remains open.

\paragraph{Other separators.}
For finite-state automata constructed from regular expressions, the implementation
obtains separators for $k > 2$ components by applying the idea of
\Cref{alg:thompson_separator} repeatedly, which adds two separator states per additional
component. With boundaries (\Cref{section:quasi_determinization_boundaries}) the growth of
$X$ matters only through the traversal, but bounds on the sizes of the resulting components
and of their boundaries remain to be proved. For general
finite-state automata with planar graphs, the separator theorem of Lipton and Tarjan
(\Cref{thm:lipton_tarjan}) gives separators of size $\mathcal{O}(\sqrt{|Q|})$. Since
the bounds depend exponentially on $|X|$, the balance between the size of the separator
and the sizes of the components determines which separators are useful in practice.
